% ===========================================================================
%  Retrieval And Subgroups
%
%  THE WRITE-UP, and it is the assistant's to type.
%    The assistant transcribes each question or topic faithfully and emits an
%    empty, marked solution region beneath it. The student does the
%    mathematics; the assistant typesets what they agreed was right, and
%    nothing that was not on the page the student sent.
% ===========================================================================
\documentclass[11pt]{article}
\usepackage{amsmath}
\usepackage{amssymb}
\usepackage[margin=1in]{geometry}
%% The two environments everything here is written in. A workspace with its own
%% `coursemacros.sty` gets that instead, and these match its definitions so a
%% document reads the same either way.
\newenvironment{problem}[1]
  {\par\medskip\noindent\textbf{Problem #1.}\ \itshape}
  {\par\medskip}
\newenvironment{solution}
  {\par\noindent\textbf{Solution.}\ }
  {\par\medskip}
\newcommand{\todo}[1]{\textbf{[TODO: #1]}}

\title{Retrieval And Subgroups}
\author{}
\date{\today}

\begin{document}
\maketitle

% One problem/solution pair per thing worked, in the order it was assigned or
% chosen. A region stays empty until the answer in it is agreed correct.

\begin{problem}{1}
Benjamini--Hochberg at $q = 0.05$ with $m = 5$ tests. The sorted p-values are
$p_{(1)},\dots,p_{(5)} = 0.004,\ 0.030,\ 0.045,\ 0.047,\ 0.060$.
The bar for rank $i$ is $\frac{i}{m}q$. Find the largest $i$ with
$p_{(i)} \le \frac{i}{m}q$ and reject ranks $1$ to $i$, where ``rejected''
means the null is rejected and the contrast is called significant.
Which ranks does BH call significant?
\end{problem}
\begin{solution}
Only the lowest rank is significant.

\begin{center}
\begin{tabular}{cccc}
rank $i$ & $p_{(i)}$ & bar $\frac{i}{5}(0.05)$ & $p_{(i)} \le$ bar? \\ \hline
1 & 0.004 & 0.01 & yes \\
2 & 0.030 & 0.02 & no \\
3 & 0.045 & 0.03 & no \\
4 & 0.047 & 0.04 & no \\
5 & 0.060 & 0.05 & no
\end{tabular}
\end{center}

The largest passing rank is $1$, so BH rejects rank $1$ alone.
\end{solution}

\end{document}
